\documentclass{article}
\usepackage{mathtools} 
\usepackage{fontspec}
\usepackage[UTF8]{ctex}
\usepackage{amsthm}
\usepackage{mdframed}
\usepackage{xcolor}
\usepackage{amssymb}
\usepackage{amsmath}


% 定义新的带灰色背景的说明环境 zremark
\newmdtheoremenv[
  backgroundcolor=gray!10,
  % 边框与背景一致，边框线会消失
  linecolor=gray!10
]{zremark}{说明}

% 通用矩阵命令: \flexmatrix{矩阵名}{元素符号}{行数}{列数}
\newcommand{\flexmatrix}[4]{
  \[
  #1 = \begin{pmatrix}
    #2_{11}     & #2_{12}     & \cdots & #2_{1#4}   \\
    #2_{21}     & #2_{22}     & \cdots & #2_{2#4}   \\
    \vdots      & \vdots      & \ddots & \vdots     \\
    #2_{#31}    & #2_{#32}    & \cdots & #2_{#3#4}
  \end{pmatrix}
  \]
}

% 简化版命令（默认矩阵名为A，元素符号为a）: \quickmatrix{行数}{列数}
\newcommand{\quickmatrix}[2]{\flexmatrix{A}{a}{#1}{#2}}

\begin{document}
\title{习题 14.2}
\author{张志聪}
\maketitle

\section*{1}

提示
\begin{align*}
  \lim\limits_{n \to \infty} \frac{a^n}{n!} = 0
\end{align*}
第二章，$\S 1$例6.

\section*{2}

\begin{itemize}
  \item (1)

        由例3可知
        \begin{align*}
          e^x = 1 + \frac{1}{1!}x + \frac{1}{2!}x^2 + \cdots + \frac{1}{n!}x^n + \cdots,
          x \in (-\infty, \infty)
        \end{align*}
        故
        \begin{align*}
          e^{x^2} = 1 + \frac{1}{1!}x^2 + \frac{1}{2!}(x^2)^2 + \cdots + \frac{1}{n!}(x^2)^n + \cdots,
          x \in (-\infty, \infty)
        \end{align*}

  \item (2)

        由例6可知（$\alpha \leq -1$）时，
        \begin{align*}
          (1 + x)^{\alpha}
           & = 1 + \alpha x + \frac{\alpha(\alpha - 1)}{2!} x^2 + \cdots + \frac{\alpha(\alpha - 1)\cdots(\alpha - n + 1)}{n!} x^n + \cdots,
          x \in (-1, 1)
        \end{align*}
        $\alpha = -1, x = -x$得到
        \begin{align*}
          \frac{x^{10}}{1 - x}
           & = x^{10} \left(1 + (-1)(-x) + \frac{(-1)(-2)}{2}(-x)^2 + \cdots + \frac{(-1)(-2)\cdots(-1 - n + 1)}{n!}(-x)^n + \cdots \right) \\
           & = x^{10} \left(1 + x + x^2 + \cdots + x^n + \cdots \right), x \in (-1, 1)                                                      \\
        \end{align*}

  \item (3)

        由例6可知（$-1 < \alpha < 0$）时，
        \begin{align*}
          (1 + x)^{\alpha}
           & = 1 + \alpha x + \frac{\alpha(\alpha - 1)}{2!} x^2 + \cdots + \frac{\alpha(\alpha - 1)\cdots(\alpha - n + 1)}{n!} x^n + \cdots,
          x \in (-1, 1]
        \end{align*}
        令$x = -2x, \alpha = -\frac{1}{2}$时得到
        \begin{align*}
          \frac{1}{\sqrt{1 - 2x}}
           & = 1 + (-\frac{1}{2}) (-2x) + \frac{(-\frac{1}{2})(-\frac{1}{2} - 1)}{2!} (-2x)^2                                                                                        \\
           & + \cdots + \frac{-\frac{1}{2}(-\frac{1}{2} - 1)\cdots(-\frac{1}{2} - n + 1)}{n!} (-2x)^n + \cdots                                                                       \\
           & = 1 + (\frac{1}{2}) (2x) + \frac{(\frac{1}{2})(\frac{1}{2} + 1)}{2!} (2x)^2 \cdots + \frac{\frac{1}{2}(\frac{1}{2} + 1)\cdots(\frac{1}{2} + n - 1)}{n!} (2x)^n + \cdots \\
           & = 1 + x + \frac{(\frac{1 \cdot 3}{2^2})}{2!} (2x)^2 \cdots + \frac{\frac{1 \cdot 3 \cdot 5 \cdots (1 + 2n - 2)}{2^n}}{n!} (2x)^n + \cdots                               \\
           & = 1 + x + \frac{1 \cdot 3}{2!} x^2 \cdots + \frac{1 \cdot 3 \cdot 5 \cdots (2n - 1)}{n!} x^n + \cdots                                                                   \\
           & = 1 + x + \frac{1 \cdot 3}{2!} x^2 \cdots + \frac{(2n - 1)!!}{n!} x^n + \cdots
        \end{align*}
        故
        \begin{align*}
          \frac{x}{\sqrt{1 - 2x}}
           & = x\left(1 + x + \frac{1 \cdot 3}{2!} x^2 \cdots + \frac{(2n - 1)!!}{n!} x^n + \cdots \right) \\
           & = x + \sum\limits_{n = 1}^{\infty} \frac{(2n - 1)!!}{n!} x^{n + 1}
        \end{align*}
        且收敛区域为
        \begin{align*}
          -1 < - 2x \leq 1 \\
          - \frac{1}{2} \leq x < \frac{1}{2}
        \end{align*}

  \item (4)

        \begin{align*}
          sin^2x = \frac{1}{2}(1 - cos2x)
        \end{align*}
        因为，在$(-\infty, \infty)$上，
        \begin{align*}
          cos x = 1 - \frac{x^2}{2!} + \frac{x^4}{4!} + \cdots + (-1)^{n} \frac{x^{2n}}{(2n)!} + \cdots
        \end{align*}
        故
        \begin{align*}
          cos 2x
           & = 1 - \frac{(2x)^2}{2!} + \frac{(2x)^4}{4!} + \cdots + (-1)^{n} \frac{(2x)^{2n}}{(2n)!} + \cdots
        \end{align*}
        所以
        \begin{align*}
          sin^2 x & = \frac{1}{2}(1 - cos2x)                                                                                                               \\
                  & = \frac{1}{2} - \frac{1}{2}\left[1 - \frac{(2x)^2}{2!} + \frac{(2x)^4}{4!} + \cdots + (-1)^{n} \frac{(2x)^{2n}}{(2n)!} + \cdots\right] \\
                  & = \frac{(2x)^2}{2!} - \frac{(2x)^4}{4!} + \cdots + (-1)^{n+1} \frac{(2x)^{2n}}{(2n)!} + \cdots
        \end{align*}

  \item (5)


\end{itemize}


\end{document}